\section{Discussion}
\label{sec:discussion}
% Course PDF: discussion = meaning, comparison, limitations.
% Points to interpret (in your words):
%  1. Why the gap grows: the defect enters the mean reference of Route A.
%     Evidence: A-clean, the contamination index, the 50% collapse.
%  2. What is NOT explained: about half of the real-defect slope remains.
%  3. The size dependence belongs to the simple mean reference, not to
%     single-image methods in general (FCA is not steeper than Route B).
%  4. Feature effect: intermediate layers raise both methods and roughly halve
%     the slope, but a significant size dependence remains at every layer on
%     the full image.
%  5. Practical meaning: what an engineer should do with large defects.
\todo{Interpretation of the findings.}

\subsection*{Limitations}
% (i) five textures, one dataset, one backbone (five layer settings, none
%     chosen on the test set); (ii) three seeds in the main table (sample SD
%     about 22% larger than the population SD shown); (iii) Route B is our
%     implementation of the AnomalyDINO idea; (iv) published numbers use other
%     protocols; (v) synthetic defects are not real defects and A-clean uses
%     the mask; (vi) border artefacts change absolute scores; (vii) Route B's
%     Pixel AP is below a published batched zero-shot result (67.3)
%     \cite{legia2025consistent}; (viii) fixed effects cannot fully separate
%     size from defect type.
\todo{Limitations.}
